\clearpage
\section{21 September 2026: Five mergers and complete land-covariate diligence}

Five documented developments now enter as one economic parent each:
Godwin/Kimberly, GO Broome, Woodside, Motto, and Elara. Their constituent
filings retain every administrative unit count, original filing date, and
refiling field. Historical parents fall from 1,814 to 1,811 and post-policy
parents from 909 to 907, while their respective housing totals remain 129,335
and 61,319. The weighting sample falls from 557 to 554 historical parents
and from 312 to 310 post-policy parents solely through these mergers.

\begin{center}
\small
\begin{tabular}{lrrr}
\hline
Development & Constituents & Housing units & Development ground \\
\hline
Godwin/Kimberly & 2 & 164 & 18,441 \\
GO Broome & 2 & 495 & 32,395 \\
Woodside & 2 & 478 & 71,862 \\
Motto & 2 & 264 & 24,880 \\
Elara & 2 & 429 & 39,003 \\
\hline
\end{tabular}
\end{center}
\noindent Ground is in square feet. These links identify economic development;
they do not establish a legal wage-assessment unit.

Motto's June 2019 environmental site plan identifies the combined development
before its December filings, including ancillary lot 155. A recorded easement
names both exact DOB jobs. The 19v1 reference parcels contain 83,936 square
feet of earlier building floor. Retained lot 37 accounts for 69,357, leaving
14,579 within the development. Elara's joint loan identifies both exact jobs,
all three parcels, and joint acquisition financing before the filings.
Woodside's six reference lots reproduce the agency's two-building site area.
GO Broome uses the common, available 19v2 reference for both filings; the
relevant parcel attributes agree with the second filing's 20v1 snapshot.

The separate inventory of 23 documented ground corrections received diligence
on earlier building area, zoning, and use. Seventeen corrections are applied:
Ocean Parkway, North 8th, East 108th, East 178th, Beach 29/30, Van Cortlandt,
Orchard, Shell Road, 21st Street, South Street, Amsterdam, West 48th, Dean,
Casa Celina, DeKalb/LIU, Crown, and Penn South. Together with the five mergers,
this batch changes 22 development records. All remain composition-eligible.
The existing parcel-allocation table now contains 34 rows for 28 developments,
including the six developments resolved before this batch.

Each applied partial-parcel correction separates included ground from earlier
buildings on retained land. At Amsterdam and Casa Celina, the retained parcel
preserves the entire earlier floor total after subdivision. At DeKalb, the
retained campus and separate garage sum exactly to the old campus floor.
At Penn South, retained campus floor plus the project's old commercial
building reproduces the earlier total. Zoning comes from the named earlier
MapPLUTO release. East 108th's frozen 2023 zoning is R5, with residential FAR
1.25 and broad FAR 2; the direct archive read corrected an erroneous R6 label
in the research note. Casa's later vacant-land classification is not backdated;
its prior-use override remains blank.

Six cases retain their previous production values. Flatbush needs an allocation
of an old building between tower and school ground. East 125th uses portions
of two earlier lots whose buildings are not located by the post-demolition
plans. Jerome's successor floor sum is 1,700 square feet below the earlier
13,000, without an explanation for that revision. The later 165th Street
survey does not locate all buildings recorded in 2023. Onderdonk's older
building has not been located confidently against the retained boundary.
Kingsbrook needs a comparable earlier floor schedule; subtracting later DOF
gross area from earlier PLUTO building area would impose an unverified
cross-source allocation. These six are queued for Jacob's judgment.

One additional conflict remains within the implemented GO Broome merger.
MapPLUTO records 4,600 square feet of earlier floor, but City Planning says
the remnants were razed in June 2019, before the reference snapshot. The
current data preserve the administrative floor value pending a decision on
using that dated agency correction. The merger and 32,395-square-foot ground
are supported independently.

The broader 126-flag inventory was held apart from the five requested merger
groups. Verification confirms that all other memberships, all constituent
unit/date/source fields, all six held parent rows, and every parent outside
the authorized batch retain their pre-batch values. Accepted areas reconcile
to the committed source table. A disposable fixture verifies fresh builds,
changed-input propagation, missing-output regeneration, and unchanged builds
on GNU Make 3.81. Reports are written alongside data and their absence does
not cause rebuilding. The older filing-map audit expected exactly 32 parents;
the mergers reduce its selected missing-map set to 29. Its checks now compare
coverage with the actual parent and map-release keys, preserving the same
matching rules. The 17 saved BIS filings now correspond to 13 parents.
Identical geometric comparisons across reference releases are counted once,
with an explicit conflict check before the parent-level join. GO Broome's
original lot list includes retained lot 1; Woodside's saved original list
covers one constituent. Exact documented crosswalks resolve those two
comparisons while preserving their raw mismatch fields.
The parent-link casebook retains its original evidence and adds the five
accepted merger records, with their primary sources and limitations. Its
current complete-characteristic sample contains 80 multi-constituent parents
(28 historical and 52 post-policy); exact evidence coverage is checked.

All five merged developments pass the refreshed boundary and filing checks.
The ten earlier flagged endpoints become five supported parents, reducing 126
flags among 869 parents to 116 among 864. Every unrelated parent's pass/flag
status is preserved. The six held floor allocations are separate from that
boundary screen, and 241 passing candidate areas still differ from production.
The wider inventory remains held for later review.

The earlier twelve-case illustrative audit also follows the merged Godwin
parent. Its September 15 geometry snapshot covered lot 79 only. A pinned
September 21 DOF response adds companion lot 88; the unchanged assessment
snapshot already contains its 5,900 square feet. Both parcels now enter the
comparison and reproduce 18,441 square feet. Original received snapshots remain
unchanged, and the source download checks its recorded hash before publication.

At the unchanged calibration specification, weights remain positive and sum
to 310; effective historical sample size is 492.89. Of 499 fixed-seed bootstrap
replications, 493 calibrate successfully and six fail explicitly. The existing
flat-charge pilot selects $\kappa=0.11$ and organization scale 1.1, with a
conditional unit reduction of 495. This is a refreshed diagnostic, not a new
specification or a claim that the unresolved land data are ready for estimation.

\textit{Sources and reproduction.} The audit note
\texttt{mergers\_and\_23\_review.md} contains the full 23-row comparison,
source-linked evidence reviews, and judgment queue. Original Motto records
and the archived site-study appendix have source hashes; automated downloads
verify those hashes before publication. Main tasks consume only the committed
manual source tables and archived administrative data. Root \texttt{make data},
\texttt{make}, \texttt{make dof-site-review}, \texttt{make pure-notch-pilot}, and
\texttt{make logbook} reproduce the affected outputs. The existing pilot is
refreshed at its unchanged specification and remains provisional while land
adjudication continues.


\clearpage
\subsection*{Flatbush follow-up: explicitly approved floor-area estimate}

Jacob approved a flagged approximation for the earlier floor on the 441-unit
Flatbush tower site. Land is 12,603 square feet, supported by the NYCIDA ground
lease and later administrative parcel area. Earlier lots 23 and 24 contribute
2,450 and 15,640 square feet of floor in full. The tower contains 61.3821\%
of old lot 18's mapped ground. Applying that fraction to its 35,580 square
feet of recorded floor gives 39,929.757 square feet of included floor and
built FAR 3.168274. The source table and canonical parent panel mark
\texttt{built\_floor\_area\_estimated = TRUE}.

The February 2018 ECF environmental review describes a three-story building
on old lot 18 and photographs its exterior. This supports a roughly uniform
height but does not establish an exact floor-space partition. The adopted
calculation assumes equal floor density across that parcel. The reason to
reject it would be a requirement for directly measured floor allocations.
Jacob selected the transparent estimate after reviewing this limitation and
its numerical consequences.

Holding corrected ground at 12,603, assigning none versus all of lot 18's
floor changes the reweighted historical share at 99 units from 1.355248\%
to 1.353476\%: a range of 0.001772 percentage points. The weighted mean ranges
from 151.3958 to 151.5662 units. The adopted allocation gives 1.354308\% and
151.4872 units. These are deliberately extreme measurement sensitivities,
not confidence intervals. The previous 1,198-square-foot parcel match yielded
151.7519 mean units. Observed housing units and counts are unchanged.

The 23-case inventory now has 18 applied corrections and five held cases.
The separate GO Broome source conflict remains open. Verification isolates
the change to Flatbush's land characteristics and the new estimate flag;
constituent records and parent memberships are unchanged. Root
\texttt{make flatbush-floor-review} reproduces the sensitivity; its research
note records source pages, the approved assumption, and the four scenarios.
